\documentclass[10pt,a4paper]{amsart}
\usepackage[margin=19mm]{geometry}
\usepackage{amsmath,amssymb,mathtools}
\usepackage[scr=rsfs,cal=euler]{mathalfa}
\usepackage{xfrac}
\usepackage[unicode,hidelinks,bookmarks=false]{hyperref}
\newcommand{\ie}{i.e.\ }
\newcommand{\cf}{cf.\ }
\newcommand{\resp}{respectively\ }
\newcommand{\Subsection}{\subsection}
\providecommand{\nobreakdash}{\nobreak\hbox{-}}
\setlength{\emergencystretch}{2em}
\multlinegap0pt
\usepackage{bm}
\usepackage{bbm}
\usepackage{dsfont}
\usepackage{xspace}
\usepackage{cancel}
\usepackage{mathtools}
\usepackage{cleveref}
\usepackage{amsthm}
\makeatletter
\@ifundefined{theorem}{\newtheorem{theorem}{Theorem}}{}
\@ifundefined{lemma}{\newtheorem{lemma}[theorem]{Lemma}}{}
\makeatother
\DeclareMathOperator{\outcomeL}{o_L}
\DeclareMathOperator{\outcomeR}{o_R}
\DeclareMathOperator{\outcome}{o}
\NewDocumentCommand\pf{smO{}}{%
  \ensuremath{
    \operatorname{pf}_{%
      \IfBooleanTF{#1}{#3}{\set{#3}}%
    }(%
      \IfBooleanTF{#1}{#2}{\set{#2}}%
    )
  }%
}
\newcommand{\set}[1]{\mathcal{#1}}
\title[On sums of $\mathscr{P}$-free forms under mis\`ere play]{On sums of $\mathscr{P}$-free forms under mis\`ere play}
\author{Alfie Davies, Sarah Miller, Rebecca Milley}
\date{Source: arXiv:2506.05257v1 (CC BY 4.0)}
\begin{document}
\maketitle

\noindent\emph{Source and licence.} On sums of $\mathscr{P}$-free forms under mis\`ere play. Version arXiv:2506.05257v1 (CC BY 4.0), distributed under the Creative Commons Attribution 4.0 International licence, as recorded in the arXiv licence metadata for this exact version and shown on the abstract page https://arxiv.org/abs/2506.05257v1. Full rights notice: licenses/arxiv-2506.05257-CC-BY-4.0.txt.\par\smallskip

\noindent\textit{Editorial scope.} Counted unit: the Lemma whose source label is \texttt{lem:B-outcome-stable} in the article (source0.tex lines 1590--1593, source label \texttt{lem:B-outcome-stable}), delivered together with the article's own complete original proof of it (source0.tex lines 1595--1623, one \texttt{proof} environment of 29 lines). No mathematics is written, repaired or re-derived: statement and proof are the article's own text line for line. Required context retained uncounted: lem:L+end (lines 1556--1560). Every definition, display and label of those lines is retained unchanged. Omitted from this package and not redistributed: 1--1555, 1561--1589, 1594--1594, 1624--2214. Nothing inside a retained excerpt is omitted or altered: every retained body is delivered exactly as the article's lines are written, so no omission receipt of any kind accompanies this package. Citations: the retained excerpts carry no citation command at all, so no bibliography entry of the article is transcribed and no reference is redistributed. Reference adaptation: none was needed and none was made; every reference inside the delivered unit resolves inside it, so no unresolved reference or citation is produced. Printed slips kept literal: no printed word, symbol, hypothesis or number of the counted statement or of its proof was altered. Layout and class: the article's own class is \texttt{article} with packages including fontenc, inputenc, babel, cfr-lm, microtype, authblk, mathrsfs, mathtools, amsfonts, amssymb, amsthm, enumerate, tikz, xcolor; the wrapper is the lane's pinned \texttt{amsart} editorial wrapper at 10pt on A4, which restates the theorem-like environments the retained text needs and adds the article's own macro definitions that the retained text needs (\texttt{\textbackslash outcome}, \texttt{\textbackslash outcomeL}, \texttt{\textbackslash outcomeR}, \texttt{\textbackslash pf}, \texttt{\textbackslash set}), with extra wrapper packages (bm, bbm, dsfont, xspace, cancel, mathtools, cleveref). The delivered numbering is the unit's own. Assets: no retained span contains a figure environment, a graphics-inclusion command, a PDF-page inclusion, a table or a TikZ picture; no image file is copied into this package, the package declares no asset dependency and no image file is redistributed with it. Licence: version arXiv:2506.05257v1 of this article is distributed under the Creative Commons Attribution 4.0 International licence, as recorded in the arXiv licence metadata for this exact version and shown on the abstract page https://arxiv.org/abs/2506.05257v1. A full rights notice is delivered at licenses/arxiv-2506.05257-CC-BY-4.0.txt. Characters: every retained excerpt is pure ASCII, so the delivered engine, XeLaTeX with its default Latin Modern text font, reports no missing character.
\begin{lemma}
    \label{lem:L+end}
    If $G,H\in\pf{B}$, where $\outcome(G)=\mathscr{L}$ and $H$ is a Left end,
    then $\outcome(G+H)=\mathscr{L}$.
\end{lemma}

\begin{lemma}
    \label{lem:B-outcome-stable}
    The blocking universe $\set{B}$ is outcome-stable.
\end{lemma}

\begin{proof}
    Let $G,H\in\pf{B}$. By symmetry, we may assume that
    $\outcome(G)=\mathscr{L}$. We need to consider two cases: when
    $\outcome(H)=\mathscr{L}$; and when $\outcome(H)=\mathscr{N}$.

    First, suppose that $\outcome(H)=\mathscr{L}$. We need to show that
    $\outcome(G+H)=\mathscr{L}$. If at least one of $G$ and $H$ is a Left end,
    then we know that $\outcome(G+H)=\mathscr{L}$ by \cref{lem:L+end}. As such,
    assume that neither $G$ nor $H$ is a Left end. Since $G$ is
    $\mathscr{P}$-free, we know that there exists some Left option $G^L$ with
    $\outcome(G^L)=\mathscr{L}$. By induction, it is clear that
    $\outcome(G^L+H)=\mathscr{L}$, and hence $\outcomeL(G+H)=\mathscr{L}$. It
    remains to show that $\outcomeR(G+H)=\mathscr{L}$. We may assume, without
    loss of generality, that Right moves to some $G+H^R$. If
    $\outcome(H^R)=\mathscr{L}$, then $\outcome(G+H^R)=\mathscr{L}$ by
    induction. Otherwise, if $\outcome(H^R)=\mathscr{N}$, then either $H^R$ is
    a Left end, in which case $\outcome(G+H^R)=\mathscr{L}$ by
    \cref{lem:L+end}, or else there exists some $H^{RL}$ with
    $\outcome(H^{RL})=\mathscr{L}$, in which case
    $\outcome(G+H^{RL})=\mathscr{L}$ by induction. Thus,
    $\outcomeR(G+H)=\mathscr{L}$.

    We now suppose that $\outcome(H)=\mathscr{N}$. We must show Left wins
    playing first on $G+H$. If $H$ is a Left end, then \cref{lem:L+end} yields
    the result. So, assume $H$ is not a Left end. Since $H$ is
    $\mathscr{P}$-free, there must exist some $H^L$ with
    $\outcome(H^L)=\mathscr{L}$. We know that $\outcome(G+H^L)=\mathscr{L}$ by
    the previous paragraph, and hence we are done.
\end{proof}

\end{document}
